% Part: first-order-logic
% Chapter: completeness
% Section: construction-of-model

\documentclass[../../../include/open-logic-section]{subfiles}

\begin{document}

\iftag{FOL}
      {\olfileid{fol}{com}{mod}}
      {\olfileid{pl}{com}{mod}}

\olsection{Konstruksi Model}

\begin{explain}
\iftag{FOL}{Saiki kita durung nggatekake $\eq$, yaiku kita mung arep
  nuduhake yen himpunan konsisten~$\Gamma$ saka !!{sentence}s kang ora
  ngemot~$\eq$ bisa dicukupi. Kaping pisan, kita ngambakake~$\Gamma$
  dadi himpunan~$\Gamma^*$ kang konsisten, !!{complete}, lan jenuh.
  Ing perkara iki, definisi model~$\Struct{M(\Gamma^*)}$ prasaja: kita
  nggunakake himpunan term katutup saka~$\Lang{L'}$ minangka domain.
  Saben !!{constant} dipasangake marang awake dhewe, lan kanthi luwih
  umum kita mesthekake yen kanggo saben term katutup~$t$,
  $\Value{t}{M(\Gamma^*)} = t$. !!{predicate}s dipasangake marang
  ekstensi saengga !!{sentence} atomik bener ing
  $\Struct{M(\Gamma^*)}$ yen lan mung yen ana ing~$\Gamma^*$. Iki
  mesthi ndadekake kabeh !!{sentence}s atomik ing~$\Gamma^*$ bener ing
  $\Struct{M(\Gamma^*)}$. Ukara liyane bener yen $\Gamma^*$ wiwitan
  iku konsisten, jangkep, lan jenuh.}{Saiki kita wis siap nemtokake
  !!a{valuation} kang ndadekake kabeh $!A \in \Gamma$ bener. Kanggo
  nindakake iki, kita dhisik ngetrapake Lemma Lindenbaum: kita entuk
  $\Gamma^* \supseteq \Gamma$ kang jangkep lan konsisten. Upamana
  !!{propositional variable}s ing~$\Gamma^*$ nemtokake~$\pAssign
  v(\Gamma^*)$.}
\end{explain}

\iftag{FOL}{%
\begin{defn}[Model term]
\ollabel{defn:termmodel}
Upamana $\Gamma^*$ minangka !!a{complete} himpunan konsisten lan jenuh
saka !!{sentence}s ing basa~$\Lang L$. \emph{Model term}~$\Struct
M(\Gamma^*)$ saka $\Gamma^*$ yaiku !!{structure} kang ditemtokake
kaya ing ngisor iki:
\begin{enumerate}
\item !!{domain}~$\Domain{M(\Gamma^*)}$ yaiku himpunan kabeh term
  katutup saka~$\Lang L$.
\item Interpretasi !!a{constant} $c$ yaiku $c$ dhewe:
  $\Assign{c}{M(\Gamma^*)} = c$.
\item !!{function}~$f$ dipasangake marang fungsi kang, nalika diwenehi
  term-term katutup $t_1$, \dots, $t_n$ minangka argumen, nduweni term
  katutup $f(t_1, \dots, t_n)$ minangka nilai:
\[
\Assign{f}{M(\Gamma^*)}(t_1, \dots, t_n) = f(t_1,\dots, t_n)
\]
\item Yen $R$ minangka !!{predicate} kanthi $n$ panggonan, mula
  \[
  \tuple{t_1, \dots,
  t_n} \in \Assign{R}{M(\Gamma^*)} \text{ yen lan mung yen } \Atom{R}{t_1, \dots,
    t_n} \in \Gamma^*.
  \]
\end{enumerate}
\end{defn}
}{%
\begin{defn}
\ollabel{defn:termmodel}
Upamana $\Gamma^*$ minangka himpunan konsisten jangkep saka
!!{formula}s. Banjur kita netepake
\[
\pAssign v(\Gamma^*)(p) = \begin{cases}
  \True & \text{yen $p \in \Gamma^*$}\\
  \False & \text{yen $p \notin \Gamma^*$}
  \end{cases}
\]
\end{defn}
}

\iftag{FOL}{%
Saiki kita bakal mriksa yen pancen $\Value{t}{M(\Gamma^*)} = t$.

\begin{lem}
\ollabel{lem:val-in-termmodel} Upamana $\Struct M(\Gamma^*)$ minangka
model term ing \olref{defn:termmodel}, banjur kanggo saben term katutup
$\Value{t}{M(\Gamma^*)} = t$. Cathetan koreksi sumber OLPL-047: bukti
lan domain model mbutuhake watesan term katutup iki.
\end{lem}
\begin{proof}
 Buktine kanthi induksi marang $t$, kanthi perkara basis nalika $t$
 minangka !!a{constant} langsung ndherek saka definisi model term.
 Kanggo langkah induksi, upamana $t_1, \ldots, t_n$ minangka term-term
 katutup saengga $\Value{t_i}{M(\Gamma^*)} = t_i$ lan $f$ minangka
 !!{function} kanthi $n$ argumen. Banjur
\begin{align*}
\Value{f(t_1,\ldots,t_n)}{M(\Gamma^*)} &= \Assign{f}{M(\Gamma^*)}(\Value{t_1}
 {M(\Gamma^*)},
\ldots, \Value{t_n}{M(\Gamma^*)}) \\
&= \Assign{f}{M(\Gamma^*)}(t_1, \dots, t_n) \\
&= f(t_1,\dots, t_n),
\end{align*}
lan mula kanthi induksi iki bener kanggo saben term katutup~$t$.
\end{proof}
}

\iftag{FOL}{%
\begin{explain}
  \iftag{prvEx}{!!^a{structure}~$\Struct{M}$ bisa ndadekake
    !!{sentence}~$\lexists[x][!A(x)]$ kang dikuantifikasi eksistensial
    bener tanpa ana conto~$!A(t)$ kang uga didadekake bener dening
    struktur kasebut.}{}
  \iftag{prvAll}{!!^a{structure}~$\Struct{M}$ bisa ndadekake kabeh
    conto~$!A(t)$ saka !!{sentence}~$\lforall[x][!A(x)]$ kang
    dikuantifikasi universal bener tanpa ndadekake
    $\lforall[x][!A(x)]$ bener.}{} Iki amarga ing umum ora saben
  !!{element} saka~$\Domain{M}$ minangka nilai term katutup
  ($\Struct{M}$ bisa uga ora katutup dening term). Iki alesan relasi
  nyukupi ditemtokake lumantar pamasangan variabel. Nanging kanggo
  model term~$\Struct{M(\Gamma^*)}$ kita, iki ora perlu amarga model
  kasebut katutup dening term. Iki isi asil sabanjure.
\end{explain}

\begin{prop}
\ollabel{prop:quant-termmodel}
Upamana $\Struct M(\Gamma^*)$ minangka model term ing
\olref{defn:termmodel}.
\begin{tagenumerate}{prvEx,prvAll}
\tagitem{prvEx}{$\Sat{M(\Gamma^*)}{\lexists[x][!A(x)]}$ yen lan mung
  yen $\Sat{M(\Gamma^*)}{!A(t)}$ kanggo paling ora siji term
  katutup~$t$.}{}
\tagitem{prvAll}{$\Sat{M(\Gamma^*)}{\lforall[x][!A(x)]}$ yen lan mung
  yen $\Sat{M(\Gamma^*)}{!A(t)}$ kanggo kabeh term katutup~$t$.}{}
\end{tagenumerate}
\end{prop}

\begin{proof}
\begin{tagenumerate}{prvEx,prvAll}
\tagitem{prvEx}{%
  \iftag{probEx}{Latihan.}{Miturut \olref[syn][ass]{prop:sat-quant},
    $\Sat{M(\Gamma^*)}{\lexists[x][!A(x)]}$ yen lan mung yen kanggo
    paling ora siji pamasangan variabel~$s$,
    $\Sat{M(\Gamma^*)}{!A(x)}[s]$. Amarga
    $\Domain{M(\Gamma^*)}$ kasusun saka term-term katutup
    ing~$\Lang{L}$, iki bener yen lan mung yen ana paling ora siji
    term katutup~$t$ saengga $s(x) = t$ lan
    $\Sat{M(\Gamma^*)}{!A(x)}[s]$. Miturut
    \olref[fol][syn][ext]{prop:ext-formulas},
    $\Sat{M(\Gamma^*)}{!A(x)}[s]$ yen lan mung yen
    $\Sat{M(\Gamma^*)}{!A(t)}[s]$, kanthi $s(x) = t$. Miturut
    \olref[fol][syn][ass]{prop:sentence-sat-true},
    $\Sat{M(\Gamma^*)}{!A(t)}[s]$ yen lan mung yen
    $\Sat{M(\Gamma^*)}{!A(t)}$, amarga $!A(t)$ iku ukara.}}{}
\tagitem{prvAll}{%
  \iftag{probAll}{Latihan.}{Miturut \olref[syn][ass]{prop:sat-quant},
    $\Sat{M(\Gamma^*)}{\lforall[x][!A(x)]}$ yen lan mung yen kanggo
    saben pamasangan variabel $s$, $\Sat{M(\Gamma^*)}{!A(x)}[s]$.
    Elinga yen $\Domain{M(\Gamma^*)}$ kasusun saka term-term katutup
    ing~$\Lang{L}$, mula kanggo saben term katutup~$t$, $s(x) = t$
    minangka pamasangan variabel kaya mangkono, lan kanggo sembarang
    pamasangan variabel, $s(x)$ minangka sawijine term katutup~$t$.
    Miturut \olref[fol][syn][ext]{prop:ext-formulas},
    $\Sat{M(\Gamma^*)}{!A(x)}[s]$ yen lan mung yen
    $\Sat{M(\Gamma^*)}{!A(t)}[s]$, kanthi $s(x) = t$. Miturut
    \olref[fol][syn][ass]{prop:sentence-sat-true},
    $\Sat{M(\Gamma^*)}{!A(t)}[s]$ yen lan mung yen
    $\Sat{M(\Gamma^*)}{!A(t)}$, amarga $!A(t)$ iku ukara.}}{}
\end{tagenumerate}
\end{proof}
}{}

\tagprob[FOL]{probEx,probAll}
\begin{prob}
Rampungna bukti \olref[fol][com][mod]{prop:quant-termmodel}.
\end{prob}
\tagendprob

\begin{lem}[Lemma Bebener]
\ollabel{lem:truth} \iftag{FOL}{Upamana $!A$ ora ngemot~$\eq$. Banjur}{}
$\iftag{FOL}{\Sat{M(\Gamma^*)}}{\pSat{v(\Gamma^*)}}{!A}$ yen lan mung
yen $!A \in \Gamma^*$.
\end{lem}

\begin{proof}
Kita mbuktekake rong arah bebarengan kanthi induksi marang $!A$.
\begin{enumerate}
\tagitem{prvFalse}{\indcase{!A}{\lfalse}
  {$\iftag{FOL}{\Sat/{M(\Gamma^*)}{\lfalse}}{\pSat/{v(\Gamma^*)}{\lfalse}}$
    miturut definisi relasi nyukupi. Ing sisih liyane, $\lfalse \notin
    \Gamma^*$ amarga $\Gamma^*$ konsisten.}}{}

\tagitem{prvTrue}{\indcase{!A}{\ltrue}
  {$\iftag{FOL}{\Sat{M(\Gamma^*)}}{\pSat{v(\Gamma^*)}}{\ltrue}$
    miturut definisi relasi nyukupi. Ing sisih liyane, $\ltrue \in
    \Gamma^*$ amarga $\Gamma^*$ konsisten lan !!{complete}, lan
    $\Gamma^* \Proves \ltrue$.}}{}

\tagitem{FOL}{\indcase{!A}{R(t_1, \dots, t_n)}
  {$\Sat{M(\Gamma^*)}{\Atom{R}{t_1, \dots, t_n}}$ yen lan mung yen
    $\tuple{t_1, \dots, t_n} \in \Assign{R}{M(\Gamma^*)}$ (miturut
    definisi relasi nyukupi), yen lan mung yen $R(t_1, \dots, t_n)
    \in \Gamma^*$ (miturut konstruksi $\Struct
    M(\Gamma^*)$).}}{\indcase{!A}{p}{$\pSat{v(\Gamma^*)}{p}$ yen lan
    mung yen $\pAssign v(\Gamma^*)(p) = \True$ (miturut definisi relasi
    nyukupi), yen lan mung yen $p \in \Gamma^*$ (miturut konstruksi
    $\pAssign v(\Gamma^*)$).}}

\tagitem{prvNot}{%
  \indcase{!A}{\lnot !B}
    {$\iftag{FOL}{\Sat{M(\Gamma^*)}}{\pSat{v(\Gamma^*)}}{\indfrm}$ yen
    lan mung yen
    $\iftag{FOL}{\Sat/{M(\Gamma^*)}{!B}}{\pSat/{v(\Gamma^*)}{!B}}$
    (miturut definisi relasi nyukupi). Miturut hipotesis induksi,
    $\iftag{FOL}{\Sat/{M(\Gamma^*)}{!B}}{\pSat/{v(\Gamma^*)}{!B}}$ yen
    lan mung yen $!B \notin \Gamma^*$. Amarga $\Gamma^*$ konsisten
    lan !!{complete}, $!B \notin \Gamma^*$ yen lan mung yen
    $\lnot !B \in \Gamma^*$.}}{}

\tagitem{prvAnd}{%
\iftag{probAnd}{%
  \indcase!{!A}{!B \land !C}{}}{%
  \indcase{!A}{!B \land
    !C}{$\iftag{FOL}{\Sat{M(\Gamma^*)}}{\pSat{v(\Gamma^*)}}{\indfrm}$
    yen lan mung yen
    $\iftag{FOL}{\Sat{M(\Gamma^*)}}{\pSat{v(\Gamma^*)}}{!B}$ lan
    $\iftag{FOL}{\Sat{M(\Gamma^*)}}{\pSat{v(\Gamma^*)}}{!C}$ kalorone
    bener (miturut definisi relasi nyukupi), yen lan mung yen
    $!B \in \Gamma^*$ lan $!C \in \Gamma^*$ kalorone (miturut
    hipotesis induksi). Miturut
    \olref[ccs]{prop:ccs}\olref[ccs]{prop:ccs-and}, iki bener yen lan
    mung yen $(!B \land !C) \in \Gamma^*$.}}}{}

\tagitem{prvOr}{%
\iftag{probOr}{%
  \indcase!{!A}{!B \lor !C}{}}{%
  \indcase{!A}{!B \lor !C}
    {$\iftag{FOL}{\Sat{M(\Gamma^*)}}{\pSat{v(\Gamma^*)}}{\indfrm}$ yen
    lan mung yen
    $\iftag{FOL}{\Sat{M(\Gamma^*)}}{\pSat{v(\Gamma^*)}}{!B}$ utawa
    $\iftag{FOL}{\Sat{M(\Gamma^*)}}{\pSat{v(\Gamma^*)}}{!C}$ (miturut
    definisi relasi nyukupi), yen lan mung yen $!B \in \Gamma^*$ utawa
    $!C \in \Gamma^*$ (miturut hipotesis induksi). Iki bener yen lan
    mung yen $(!B \lor !C) \in \Gamma^*$ (miturut
    \olref[ccs]{prop:ccs}\olref[ccs]{prop:ccs-or}).}}}{}

\tagitem{prvIf}{%
\iftag{probIf}{%
  \indcase!{!A}{!B \lif !C}{}}{%
  \indcase{!A}{!B \lif !C}
    {$\iftag{FOL}{\Sat{M(\Gamma^*)}}{\pSat{v(\Gamma^*)}}{\indfrm}$ yen
    lan mung yen
    $\iftag{FOL}{\Sat/{M(\Gamma^*)}}{\pSat/{v(\Gamma^*)}}{!B}$ utawa
    $\iftag{FOL}{\Sat{M (\Gamma^*)}}{\pSat{v(\Gamma^*)}}{!C}$ (miturut
    definisi relasi nyukupi), yen lan mung yen $!B \notin \Gamma^*$
    utawa $!C \in \Gamma^*$ (miturut hipotesis induksi). Iki bener yen
    lan mung yen $(!B \lif !C) \in \Gamma^*$ (miturut
    \olref[ccs]{prop:ccs}\olref[ccs]{prop:ccs-if}).}}}{}

\iftag{FOL}{%
\tagitem{prvAll}{%
  \iftag{probAll}{%
    \indcase!{!A}{\lforall[x][!B(x)]}{}}{%
    \indcase{!A}{\lforall[x][!B(x)]}{$\Sat{M(\Gamma^*)}{\indfrm}$ yen
      lan mung yen $\Sat{M(\Gamma^*)}{!B(t)}$ kanggo kabeh term
      katutup~$t$ (\olref{prop:quant-termmodel}). Miturut hipotesis
      induksi, iki bener yen lan mung yen $!B(t) \in \Gamma^*$ kanggo
      kabeh term katutup~$t$; miturut
      \olref[hen]{prop:saturated-instances}, iki sabanjure bener yen lan
      mung yen $\lforall[x][!B(x)] \in \Gamma^*$. Cathetan koreksi
      sumber OLPL-048: rong kuantifikasi term iki diwatesi marang term
      katutup, kaya pratelan kang dirujuk. OLPL-049: formula pungkasan
      nggunakake subformula B, dudu formula induksi njaba.}}}{}

\tagitem{prvEx}{%
  \iftag{probEx}{%
    \indcase!{!A}{\lexists[x][!B(x)]}{}}{%
    \indcase{!A}{\lexists[x][!B(x)]}{$\Sat{M(\Gamma^*)}{\indfrm}$ yen
      lan mung yen $\Sat{M(\Gamma^*)}{!B(t)}$ kanggo paling ora siji
      term katutup~$t$ (\olref{prop:quant-termmodel}). Miturut hipotesis
      induksi, iki bener yen lan mung yen $!B(t) \in \Gamma^*$ kanggo
      paling ora siji term katutup~$t$. Miturut
      \olref[hen]{prop:saturated-instances}, iki sabanjure bener yen lan
      mung yen $\lexists[x][!B(x)] \in \Gamma^*$. Cathetan koreksi
      sumber OLPL-048: loro-lorone nggunakake term katutup.}}}{}
}{}
\end{enumerate}
\end{proof}


\tagprob[FOL]{probOr,probAnd,probIf,probEx,probAll}

\begin{prob}
Rampungna bukti \olref[fol][com][mod]{lem:truth}.
\end{prob}

\tagendprob

\tagprob[notFOL]{probOr,probAnd,probIf}

\begin{prob}
Rampungna bukti \olref[pl][com][mod]{lem:truth}.
\end{prob}

\tagendprob

\end{document}
